\begin{theorem}[Source result P3K.109 under continuation review]\label{night:P3K.109}
\StatementRecord{night:P3K.109}{open}{}{\cref{night:P3K.94}, \cref{night:P3K.96}, \cref{night:P3F.35}, \cref{night:P3F.46}; the original source, point line, roots, characters, reference, selected fan and scalar range stated here}
Set
\[
 P_r=\frac{324r^3-405r^2+144r-16}{2187},\qquad
 Q_r=\frac{2(9r-2)(972r^3-891r^2+288r-32)}{531441},
\]
and let \(r\) lie in the explicit source locus
\[
 \mathcal U=\{r\in\mathbb C:r(4r-1)(3r-1)P_rQ_r\ne0\}.
\]
Retain the original source
\[
 E_r:\ y^2=x^3+rx^2+tx+t^2,\qquad G=(0,t),
\]
one chosen root \(t_c\) of
\(27t^2+(4-18r)t+r^2(4r-1)=0\), the selected height
\(q=t-t_c\), its actual branch marking, and any prescribed
lattice-basis fan. Let the point line and suspension be
\(M_a=\mathcal O(aG-O)\) and translation by \(-bG\), with
either of the complete cap progressions
\[
 (a,b;m)=(6,1;3l+2),\quad l\ge0,
 \qquad
 (a,b;m)=(2,3;3l),\quad l\ge1.
\]
Use either signed primitive norm \(n=\pm1\), the prescribed finite
character \(\zeta_m^d\), the full annular reference of P3K.96,
the ordinary attachments at the other finite source fibres, and a
fixed nonzero scalar in the strict contraction range. Let \(X\)
be the resulting completed complex threefold.

Multiplication on the genuine point line, carried through the full
root, reference and toric maps, gives the exact group isomorphism
\[
       \mathbb C^*_{\rm physical}
          \xrightarrow{\ \cong\ }\operatorname{Aut}_{\rm hol}(X).
\]
Every holomorphic automorphism preserves the original base pointwise
and is isotopic to the identity through this physical action. For
two constructed holomorphic endpoint identifications of the same
cut family, their ratio is precisely such a physical scalar; the
inverse-endpoint returns are therefore isotopic by the displayed
completed maps.

The proof includes all possible descended base maps on this explicit
source locus, both cap progressions, both norm signs and every
retained finite character. It assigns no new source or scalar return
and no shorter endpoint outside the constructed family. A concrete
object is
\[
 r=\tfrac12,\quad t_c=(5+i\sqrt2)/54,\quad
 (a,b;m,n,d)=(6,1;2,1,0),\quad T_0=T(e_1,e_2),
\]
with the complete P3K.94--96 reference.
\end{theorem}
\begin{proof}[Supplied source argument; continuation review in progress]

% Expanded verified cp36 input: notes/part3_e6_fixed_source_endpoints_proof.tex
\par

\medskip\noindent\textbf{EF.1. The actual source locus, line, cap and first object}

\par\nobreak

Put
\[
 E_r:\quad y^2=x^3+rx^2+tx+t^2,\qquad G=(0,t),\qquad
 s=\frac14-r.
 \tag{EF.1}
\]
Write
\[
 \begin{split}
 S_r&=(18r-4)/27,\quad C_r=r^2(4r-1)/27,\\
 P_r&=(324r^3-405r^2+144r-16)/2187,\\
 Q_r&=2(9r-2)(972r^3-891r^2+288r-32)/531441.
 \end{split}                                                    \tag{EF.2}
\]
The source locus in this proof is the explicit open set
\[
 \mathcal U=\{r\in\mathbb C:
 r(4r-1)(3r-1)P_rQ_r\ne0\}.
 \tag{EF.3}
\]
Choose one of the two roots of
\(27t^2+(4-18r)t+r^2(4r-1)=0\), call it \(t_c\), and retain
that marked selected I1, its actual branch coordinate and its prescribed
lattice-basis fan. Write \(q=t-t_c\). The other I1 and the I2 at zero
have their ordinary source attachments. There is one finite cap at the
original infinity, with the entire reference map of K96.

The two data in the scope are
\[
 (a,b;m)=(6,1;3l+2),\ l\ge0,
 \qquad (a,b;m)=(2,3;3l),\ l\ge1,
 \qquad n\in\{1,-1\}.
 \tag{EF.4}
\]
The point line is the actual \(M_a=\mathcal O(aG-O)\); the source
suspension is translation by \(-bG\). The character \(\zeta_m^d\)
and the whole annular reference cochain are retained. The scalar is a
fixed nonzero value in the contraction domain of the construction. The
source, root and line maps in EX.1-\/-19 of K96 construct these actual
models. The integral maps of F46 and the common blowup, two-ruling,
Gysin and van Kampen maps of F35 give the sphere on each prescribed fan.
Equivalently, one can transport the explicit K96 sphere along a compact
path of marked simple selected nodes to the chosen point of EF.3: all
EX charts are holomorphic along that path, and a common smaller scalar
gives its strict compact norm shell. Changing that scalar within the
contraction domain gives the same smooth sphere. Thus the completed
threefold used here has \(H^2=H^4=0\), with its actual source fibration.

A fully specified first object is
\[
 r=\tfrac12,\quad t_c=(5+i\sqrt2)/54,\quad
 (a,b;m,n,d)=(6,1;2,1,0),\quad T_0=T(e_1,e_2),
 \tag{EF.5}
\]
with the fixed small scalar and the K94/K96 reference lift. Here
\(P_r=-19/8748\) and \(Q_r=215/2125764\), so it belongs to EF.3.

The conclusion proved below is the exact group equality
\[
       \operatorname{Aut}_{\rm hol}(X)=\mathbb C^*_{\rm physical}.
 \tag{EF.6}
\]
The right-hand map is multiplication on the actual point line, followed
through all root, annular and Cox arrows. In particular it is injective
and every endpoint automorphism is joined to the identity by that action.

\par

\medskip\noindent\textbf{EF.2. Recover the base and fix it by the actual critical values}

\par\nobreak

Let \(\pi:X\to\mathbb P^1_t\) be the constructed fibration. A smooth
fibre is a compact complex two-torus. Give it any flat Kähler form
\(\omega\). A horizontal effective divisor on \(X\) would restrict
to a nonzero effective divisor \(D_T\) on a generic fibre. Since its
line class restricts from \(H^2(X)=0\),
\[
 0=\int_T c_1(\mathcal O(D_T))\wedge\omega
   =\int_{D_T}\omega>0.
 \tag{EF.7}
\]
The divisor fundamental-class identity includes every multiplicity, so
this excludes every horizontal polar divisor of a meromorphic function.
The function is constant on each generic compact fibre, by removal of
its remaining codimension-two indeterminacy there. It therefore descends
on the regular base. At a generic smooth point of a reduced special-fibre
component choose a transverse disc. After taking the extending root of
its nonzero base unit, its exact base map is \(t-t_0=u^e\).
The pulled-back meromorphic function has a Laurent expansion with finite
principal part. Constancy on the generic fibres makes it invariant under
\(u\mapsto\zeta_eu\), so every occurring exponent is a multiple of
\(e\). This gives its meromorphic base extension at \(t_0\). Thus
\(\pi^*\) identifies \(\mathbb C(X)\) with \(\mathbb C(t)\).
Every biholomorphism consequently has the exact equation
\(\pi F=h\pi\) for a Möbius map \(h\), first on the regular locus
and then on the whole space.

The direct source discriminant is
\[
 -t^2\{27t^2+(4-18r)t+r^2(4r-1)\},
 \qquad \operatorname{disc}_t\{\cdots\}=16(1-3r)^3.
 \tag{EF.8}
\]
The local actual source maps at the two simple nodes are \(xy=q_i\)
after the nonzero analytic Morse factor; their invariant differential
has an annular logarithmic period. A positive node meridian therefore
has a nontrivial unipotent elliptic block. The I2 has its two-node
unipotent block. The selected toric characters retain the same base
function \(q=t-t_c\), and introduce no other base critical value.
The ordinary source/line quotients are submersions at every other finite
value. Thus the finite critical set is exactly \(\{0,t_+,t_-\}\).

At infinity use the actual minimal good cubic source root
\[
 q_\infty=z^3,\quad x_0=z^4X,\quad y_0=z^6Y,\qquad
 Y^2=X^3+r z^2X^2+zX+1.
 \tag{EF.9}
\]
Its deck is \((z,X,Y)\mapsto(\zeta_3z,\zeta_3^2X,Y)\).
On the central curve the differential has eigenvalue \(\zeta_3^2\),
so the elliptic monodromy polynomial is \(x^2+x+1\). The full cap
annulus of EF.7 below conjugates the original periods with their
functions, and retains this elliptic quotient. Its full rank-four
polynomial is \((x-1)^2(x^2+x+1)\). Infinity is therefore the unique
critical value with these non-unipotent eigenvalues. The base map fixes
infinity and is affine on \(\mathbb C_t\).

It permutes the finite set in EF.8. Its centroid is \(S_r/3\).
In the centered coordinate \(u=t-S_r/3\), the monic polynomial of
that set is exactly \(u^3+P_ru+Q_r\). An affine permutation fixes
the centroid and has the form \(u\mapsto c u\). Comparing its two
nonzero coefficients gives \(c^2=c^3=1\), hence \(c=1\).
This proves \(h=1\) on the stated source locus, without presupposing
that the automorphism preserves the elliptic source or its marked point.

\par

\medskip\noindent\textbf{EF.3. The actual point-theta period module}

\par\nobreak

On a simply connected regular base chart choose oriented elliptic
periods, set \(\tau=\omega_2/\omega_1\), and let \(g\) be the
chosen Abel integral of \(G\), divided by \(\omega_1\). Put
\(Q=e^{2\pi i\tau}\), \(P=e^{2\pi ig}\), and
\(Z=e^{2\pi iz}\). Write
\[
 f_R(Z)=\frac{\Theta(Z/R;Q)}{\Theta(Z;Q)},\qquad
 V_{\rm pt}=f_{P^a}(Z)v_M.
 \tag{EF.10}
\]
This is the coefficient of the actual point line in its constant-factor
frame. In additive coordinates \((z,w)\), with
\(V_{\rm pt}=e^{2\pi iw}\), the four literal periods, ordered as
\((\gamma,e_1,e_2,\delta)\), are
\[
 \gamma=(-bg,\ell),\qquad e_1=(1,0),\qquad
 e_2=(\tau,ag),\qquad \delta=(0,1).
 \tag{EF.11}
\]
In this frame the complete suspension coefficient is constant on each
smooth elliptic fibre: the ratio with the translation lift having the
same constant factors is an everywhere nonzero holomorphic function on
that compact curve. This proves the constant-factor assertion while
retaining its actual base value. Here \(\ell\) is the chosen logarithm
of the \textbf{entire} multiplier
of \(B_\lambda=\lambda B_{\rm sel}\) in the point frame. It includes
\(q/t^3\), the full Miller coefficient, the original scalar and every
Cartier unit. It is never replaced by a leading coefficient.

For comparison with the earlier tensor coordinates retain the exact
typed identity
\[
 V_{\rm tensor}=f_P(Z)^a f_a(p)^{-1}v_M
       =C_a(t)V_{\rm pt},\qquad
 C_a(t)=\frac{f_P(Z)^a}{f_a(p)f_{P^a}(Z)}.
 \tag{EF.12}
\]
The full Miller divisor identity makes this ratio independent of the
source point. Its complete base value and its continuation are retained
in \(\ell\). EF.12 is used in both directions at each original
Cartier chart, including coincident sections.

There are two exact changes in EF.11. First, an Abel representative
change \(g'=g+\mu+\nu\tau\), with integral \(\mu,\nu\), gives
\(P'^a=P^aQ^{a\nu}\). The identity
\(\Theta(Q^{-k}z)=(-z)^kQ^{-k(k+1)/2}\Theta(z)\) gives the full arrow
\[
 V_{\rm pt}'=(-Z/P^a)^{a\nu}
            Q^{-a\nu(a\nu+1)/2}V_{\rm pt}.
 \tag{EF.13}
\]
Its inverse uses the reciprocal at the same source point. Pulling the
new periods through that inverse gives
\[
 e_1'=e_1-a\nu\delta,\qquad
 e_2'=e_2+a\mu\delta,\qquad
 \gamma'=\gamma-b\mu e_1-b\nu e_2+c\delta,
 \tag{EF.14}
\]
where the complete integer is
\[
 c=\ell'-\ell+ab\nu(2g+\mu+\nu\tau).
 \tag{EF.15}
\]
This is an integer because these are periods of the same actual closed
torus map. Formula EF.15 specifies it using the entire logarithms; it
does not set it to zero. All signs and the source-period reindex are
fixed by the inverse of EF.13.

Second, for a change of oriented elliptic basis
\(A=\left(\begin{smallmatrix}p&q_0\\r_0&s_0\end{smallmatrix}\right)
\in SL_2(\mathbb Z)\), put \(\lambda_A=p+r_0\tau\). The directed
source and line map is
\[
 z'=z/\lambda_A,\quad \tau'=(q_0+s_0\tau)/\lambda_A,
 \quad g'=g/\lambda_A,\quad
 w'=w-\frac{a r_0g}{\lambda_A}z+c_A.
 \tag{EF.16}
\]
The full constant is defined by the actual point-divisor arrow:
\[
 e^{2\pi i c_A}=
 \frac{f_{e^{2\pi iag'}}(e^{2\pi iz'};e^{2\pi i\tau'})}
      {f_{P^a}(Z;Q)}
 \exp\!\left(2\pi i\frac{a r_0g}{\lambda_A}z\right).
 \tag{EF.17}
\]
The numerator and denominator have the same transported point divisor.
Their period factors cancel by EF.16, so the right side is a nonzero
constant on the compact elliptic fibre. This proves the displayed
Cartier arrow, and specifies its complete value by evaluation at any
ordinary point. Its logarithm branch is carried into \(\ell\).
The two old vectors \(p e_1+r_0e_2\), \(q_0e_1+s_0e_2\) go
exactly to \((1,0),(\tau',ag')\). The suspension line coefficient
changes by \(ab r_0g^2/\lambda_A\), before its retained integral
logarithm adjustment. Thus this change introduces no unrecorded mixed
period term.

Combining these two actual arrows gives the monodromy on the full
integral lattice:
\[
 M_\alpha=
 \begin{pmatrix}
 1&0&0\\
 b v_\alpha&A_\alpha&0\\
 c_\alpha&a v_\alpha^t J A_\alpha&1
 \end{pmatrix},\qquad
 J=\begin{pmatrix}0&-1\\1&0\end{pmatrix}.
 \tag{EF.18}
\]
The first and last block sizes are one. In EF.14,
\(v=(-\mu,-\nu)^t\). In the complete atlas the cocycle and its
central cochain obey the literal product equations
\[
 v_{\alpha\beta}=v_\alpha+A_\alpha v_\beta,\qquad
 c_{\alpha\beta}=c_\alpha+c_\beta+
             ab v_\alpha^t J A_\alpha v_\beta.
 \tag{EF.19}
\]
They are the multiplication equations of the displayed integral
matrices. The line scalar loop \(\delta\) is primitive, as is the
quotient homomorphism \(\psi(\gamma)=1\),
\(\psi(e_1)=\psi(e_2)=\psi(\delta)=0\): the nonzero-line cover
has fundamental fibre lattice generated by \(e_1,e_2,\delta\),
and adjoining the actual suspension adds its primitive deck generator.
EF.7 below retains these generators on the full cap annulus.

\par

\medskip\noindent\textbf{EF.4. The Abel class is nonzero and the full centralizer is computed}

\par\nobreak

The cocycle \(v_\alpha\) of EF.18 is nonzero in rational cohomology.
Here is a direct proof on the actual source. A rational coboundary would,
after multiplication by an integer and subtraction of that integral
period cochain, give a single-valued holomorphic Lie logarithm of
\(NG\) on the regular base. Multiply that integer further by six.
Then the section lies in the identity component at every special fibre.

At each I1 and at I2 the actual Tate identity coordinate of \(NG\)
is an extending nonzero unit. Its chosen analytic logarithm gives a
local extending Lie section. The difference from the purported global
logarithm is a monodromy-invariant period. For an I1 or I2 the invariant
period is the annular vanishing-cycle period; integrating the invariant
differential in \(xy=q_i\) shows that it is a holomorphic unit times
\(2\pi i\). It therefore extends in the same Lie frame. There is
no logarithmic second period in an invariant vector.

At infinity the actual formulas are particularly explicit. The chords
give
\[
 \begin{split}
 2G&=(s,-t-s/2),\\
 x(3G)&=4t^2/s^2+2t/s,\\
 y(3G)&=8t^3/s^3+6t^2/s^2+t/s-t.
 \end{split}                                                    \tag{EF.20}
\]
In EF.9 the root formal logarithm of \(3G'\) is
\(-s z^2/2+O(z^5)\). The differential in the original minimal infinity chart is
\(\omega_0=dx_0/(2y_0)=z^{-2}\omega_{\rm good}\); the affine
differential is instead \(\omega_t=z\omega_{\rm good}\). The
root logarithm transforms by \(\zeta_3^2\) under the actual deck,
so its quotient by \(z^2\) is invariant and descends holomorphically
in \(q_\infty=z^3\). Consequently division
of that root logarithm by \(z^2\) is the extending original Lie
logarithm. Integer multiples retain this extension. The difference
between the local and purported global logarithms is an invariant
elliptic period, and EF.9 has no nonzero invariant period. Thus this
difference is zero, including its full branch.

The resulting logarithm is a global holomorphic section of the actual
Lie line. In the original coordinates
\(x_0=q_\infty^2x,y_0=q_\infty^3y\), its differential frames obey
\(\omega_0=q_\infty^{-1}\omega_t\), so its dual frames obey
\(W_0=q_\infty W_t\). The Lie line is therefore \(\mathcal O(-1)\).
In coefficients, an entire function \(u_t(t)\) must satisfy
\(u_0(q_\infty)=q_\infty^{-1}u_t(q_\infty^{-1})\) holomorphic
at zero. Hence \(u_t(t)=O(t^{-1})\) at infinity, and is zero.
This would give \(NG=O\).

The explicit original central coordinate disproves that conclusion.
In the IV* Weierstrass chart EF.20 gives
\[
 x_0(3G)=4/s^2,\qquad y_0(3G)=8/s^3,\qquad
                 -x_0(3G)/y_0(3G)=-s/2\ne0.
 \tag{EF.21}
\]
The parameter \(-x_0/y_0\) is the additive coordinate on the smooth
identity component of the central cusp; direct substitution
\((x_0,y_0)=(u^2,u^3)\) gives
\(dx_0/(2y_0)=d(-1/u)\). Its integer multiples are nonzero.
Thus \(G\) is not torsion and the rational Abel cocycle is nonzero.

The elliptic monodromy centralizer is \(\mathbb Q I\). Indeed the
infinity matrix has irreducible polynomial \(x^2+x+1\), so its
rational centralizer is the quadratic field it generates. A nontrivial
I1 unipotent cannot belong to this field: its nonzero nilpotent part
would be a zero divisor. A non-scalar element of that quadratic field
generates it, so no non-scalar element commutes with the I1 matrix.

We now solve the full rank-four problem, including every
\(c_\alpha\). Modulo \(\delta\), an invariant vector with
\(\gamma\) coefficient \(z\) would give
\((A_\alpha-I)u+bzv_\alpha=0\). The nonzero Abel class forces
\(z=0\), and the absence of elliptic invariants forces \(u=0\).
Thus the invariant line of the full module is exactly
\(\mathbb Q\delta\). The dual computation using
\(v_\alpha^tJA_\alpha\) gives the invariant covector line exactly
\(\mathbb Q\psi\). A commuting matrix must preserve both and has
the form
\[
 X=\begin{pmatrix}
 s_0&0&0\\ u&A&0\\ c&w&t_0
 \end{pmatrix}.
 \tag{EF.22}
\]
Its elliptic quotient is \(A=eI\). The middle-left commutation
equation is
\((I-A_\alpha)u=b(s_0-e)v_\alpha\), hence \(s_0=e,u=0\).
The lower-middle equation is
\(w(A_\alpha-I)=a(e-t_0)v_\alpha^tJA_\alpha\).
Writing \(w=z^tJ\) and using \(A_\alpha^tJA_\alpha=J\) turns
this into the same Abel coboundary equation. Therefore
\(t_0=e,w=0\). The remaining equations allow precisely
\[
                 X=eI+c\delta\psi.
 \tag{EF.23}
\]
They include the full central cochain, because this matrix commutes
with every EF.18 for its actual value of \(c_\alpha\). Integral
unimodularity gives \(e=\pm1\), \(c\in\mathbb Z\).

A fibre biholomorphism is affine: its lifted derivative is periodic,
so its entries are holomorphic functions on the compact torus and are
constant. Its real linear map is complex linear. Subtracting
\(eI\) from EF.23 leaves the real rank-one map \(c\delta\psi\)
when \(c\ne0\); a complex-linear nonzero map has even real rank.
Consequently every actual holomorphic endpoint has linear part exactly
\(+I\) or \(-I\).

\par

\medskip\noindent\textbf{EF.5. The completed fixed-curve obstruction excludes the negative part}

\par\nobreak

Take the actual cap-completed automorphism with linear part \(-I\)
whose possibility remains in EF.4. On each regular two-torus it has
the affine form \(x\mapsto-x+a(t)\). Multiplication by two on its
integral lattice has degree \(2^4=16\), so the actual fixed set over
the regular base has exactly sixteen points in each fibre. Its vertical
derivative is \(-2I\); the analytic implicit-function theorem makes
this fixed set a finite étale curve there.

The global equalizer of the two actual holomorphic maps \(F\) and
identity is a closed analytic subset of the compact completed \(X\).
Choose an irreducible component containing one of these regular fixed
points. Its local dimension is one, hence its entire irreducible
component is a compact analytic curve. This is the proper analytic
closure of that horizontal branch, with all source, toric and cap
limits retained. Its map to the base has positive degree. Equivalently,
its intersection with a generic smooth torus fibre is a positive sum
of the transverse complex intersection multiplicities.

But \(H^4(X;\mathbb Z)=0\), so the Poincaré dual of this curve is
zero and the same generic intersection is zero. One may also pair its
normalization with the exact class \(\pi^*[\omega_{\mathbb P^1}]\)
in \(H^2(X)=0\): its positive degree gives a positive integral and
Stokes gives zero. This contradiction concerns the actual global
completed map. It does not use a local inverse whose cap extension was
only postulated. Thus every holomorphic endpoint has linear part \(+I\).

\par

\medskip\noindent\textbf{EF.6. The identity-linear endpoint factors through the actual scalar}

\par\nobreak

On the regular fibres an identity-linear map is an actual translation.
It preserves the primitive \(\psi\) and lifts locally to the cyclic
nonzero-line cover. It commutes with its generator \(B_\lambda\).
Its lifted affine map is equivariant for the full physical \(\mathbb
C^*\), and its source projection is an elliptic translation.

At an ordinary critical disc, this cover is the actual quotient of
\(M_a^\times\) over the resolved source by the invertible contraction
\(B_\lambda\). The base disc is simply connected and its total
fundamental group is obtained from the generic fibre by its vanishing
cycles. Those cycles lie in the kernel of \(\psi\), as the displayed
cover shows. Thus the preserved fibre homomorphism extends and the map
lifts holomorphically across that disc. Equivariance and source
projection extend from the dense regular locus. Evaluation on the
chosen nonzero vector over \(O\) gives the actual extending source
section there.

At infinity take the literal fibre product with \(z^m=q_\infty\)
and its normalization. EF.7 identifies it with the root quotient by
\(B_\lambda\), because the kernel of the full cap generator has
\(A_\lambda^m=B_\lambda^n\), \(n=\pm1\). The canonical lift of
the base-preserving map fixes \(z\). Its nonzero-line lift has a
commutator \(B_\lambda^k\) with the actual \(A_\lambda\).
Iteration \(m\) times, and its already proved commutation with
\(B_\lambda^n\), give \(B_\lambda^{mk}=1\). Strict contraction
gives infinite order, hence \(k=0\). This retains the entire finite
character and both signed norms. Source evaluation therefore extends
on the genuine root and on its source graph, not merely on a generic
elliptic fibre.

At the selected I1 choose a source section in an open toric stratum.
In a chart with the actual equation \(q=z_0z_1z_2\), its image under
the actual endpoint is a section over \(q\). The vanishing orders of
these three holomorphic coordinates are nonnegative and sum to one;
the source and line torus characters consequently have integral finite
orders. Write the complete periods there as
\[
 h:(Z,V,q)\mapsto(QZ,\alpha V,q),\qquad
 B:(Z,V,q)\mapsto(gZ,q\eta V,q),\quad Q=\kappa q,
 \tag{EF.24}
\]
where \(\kappa,\alpha,g,\eta\) are the actual nonzero holomorphic
source/Cartier factors. Reindex by the literal \(h^rB^k\) to remove
the two integral orders; retain its full unit factors. Evaluation of
the endpoint and its inverse then gives extending nonzero units
\(K(q),C(q)\). On every fan character their exact map is
\[
 Z^uV^vq^h\longmapsto K(q)^uC(q)^vZ^uV^vq^h.
 \tag{EF.25}
\]
The inverses use the inverse units at the actual inverse point. This
proves source-section extension across the selected I1 and the full
bounded line arrow on every prescribed cone; the integral \(B\)
reindex will be used below, not deleted.

Away from the selected point the local lifts differ by integral
suspension powers. The orbifold base is \(\mathbb P^1\setminus
\{t_c\}\) with the retained order-\(m\) cap; its fundamental group
is the finite cyclic cap group. A homomorphism to \(\mathbb Z\) is
zero. The integer cocycle of lifts is therefore removed by its actual
deck powers. The resulting source section is global on the minimal
elliptic surface.

For clarity the marked source generator is identified integrally.
The explicit rational chart
\[
 U=y-t-x/2,\quad V=y+t+x/2,\quad UV=x^2(x-s)
 \tag{EF.26}
\]
shows the source is rational. Its resolved elliptic surface has Euler
number \(8+2+1+1=12\), hence the unimodular rank-ten rational-surface
intersection lattice. The fibre and zero section give the unimodular
rank-two summand. The actual IV* and I2 component lattices are E6 and
A1, with determinant six. Restriction of a divisor to the generic
fibre identifies the quotient by these vertical classes with the
Mordell-\/-Weil group, of free rank one. The actual component and
intersection arrows can be read directly. At finite base values \(G\)
is in the affine chart and \(O\) is at infinity. At the A1 point of
EF.26 the tangent cone is \(UV=-s x^2\). The section \(G\) has direction
\([x:U:V]=[0:0:1]\), whereas the strict central fibre has equation
\(V-U-x=0\). Thus the lift of \(G\) meets the exceptional nonidentity
A1 component once, off its intersection with the identity component.
The two sections remain disjoint. At infinity, EF.9 has the distinct
fixed sections \((0,1)\) and \(O\) on the good central elliptic curve.
The actual quotient charts at each have weights \((2,1)\), resolved
by \((2/3,1/3),(1/3,2/3)\), so their section arms remain disjoint.
Label the nonidentity outer component met by \(G\) as \(C_0\), and
order the other E6 components by the edges
\((0,1),(1,2),(2,3),(3,4),(2,5)\). The arm through \(O\) has outer
component \(F-C_0-2C_1-3C_2-2C_3-C_4-2C_5\).
These are the literal component labels of F41 SC.5 and the integral
source matrix in J.4, carried by the displayed root charts.

The positive Cartan matrix and actual correction column are
\[
 A=\begin{pmatrix}
 2&-1&0&0&0&0\\-1&2&-1&0&0&0\\
 0&-1&2&-1&0&-1\\0&0&-1&2&-1&0\\
 0&0&0&-1&2&0\\0&0&-1&0&0&2
 \end{pmatrix},\quad
 A^{-1}e_0=(4/3,5/3,2,4/3,2/3,1)^t.
 \tag{EF.26a}
\]
The section intersection vector is exactly \(e_0\); at A1 it is
\((1)\), with Cartan matrix \((2)\). Hence the two local corrections
are \(e_0^tA^{-1}e_0=4/3\) and \(1/2\). The displayed disjoint
section charts prove \(G\cdot O=0\), so the height is
\(2-4/3-1/2=1/6\). The normal line of either smooth section is the
actual Lie line, so \(O^2=G^2=-1\). In the order
\((O,F,C_0,\ldots,C_5,A_1,G)\), put \(OF=GF=1\), the component block
\(-A\oplus(-2)\), \(G C_0=G A_1=1\), and all other unspecified
pairings zero. This full integral Gram matrix has determinant \(-1\),
as checked directly. The ten actual classes therefore give the whole
unimodular source lattice, and its quotient by the first nine is the
free group generated by \(G\). The height determinant gives the same
primitive conclusion explicitly. Writing \(\iota\) for its free index and
\(t_{\rm MW}\) for the torsion order, the integral determinant
equations are
\[
 1=6h(G_{\rm prim})/t_{\rm MW}^2,\qquad
 \iota^2h(G_{\rm prim})=1/6,\qquad
                         \iota^2t_{\rm MW}^2=1.
 \tag{EF.27}
\]
Both positive integers are one. Thus the actual section above is
\(jG\) for an integer \(j\), with no additional torsion section.

Keep the full Miller functions of K96. Their typed one-step map is
\[
 B_{a,1}(p)=f_2(p)^a\frac{f_a(p-G)}{f_a(p)},\quad
 U_G=B_{a,1}^{-1},\quad
 U_G(p,v)=(p+G,B_{a,1}(p+G)^{-1}v).
 \tag{EF.28}
\]
All inverses therefore evaluate at the actual target source point.
The chronological product \(B_{a,b}\) is \(U_G^{-b}\), and
\[
 B_{\rm sel}=(q/t^3)B_{a,b},\qquad
              B_{\rm sel}^{-1}=(t^3/q)U_G^b.
 \tag{EF.29}
\]
An endpoint lift over \(+jG\) is \(H(t)U_G^j\): the ratio is an
invertible function on each compact elliptic fibre and is a base
function. Its meromorphic extension at the finite singular fibres and
root is in the actual Cartier module charts. The selected argument
EF.24 gives its finite integer order there. Hence \(H\) is a rational
base function.

The actual generic-component valuations are
\[
\begin{array}{c|ccc}
 &t=0&z=0\ (q_\infty=z^m)&q=0\\ \hline
 U_G&-3/b&2m/b&0\\
 H&3j/b&-2mj/b&-k
\end{array}                                                     \tag{EF.30}
\]
Here are the exact full-arrow antecedents of these orders. For M2,
EX.1 gives \(B_{2,1}(p)=f_2(p)f_2(p-G)\). On the identity I2
component \(f_2\) has order zero; on its other component the actual
chart \(x=t\alpha,y=t\beta\) gives
\[
 f_2=t\,\frac{\beta-1-\alpha/2}{-s+t\alpha}.
 \tag{EF.30a}
\]
Both displayed generic coefficients are nonzero. The actual \(-G\)
map exchanges the components, so \(B_{2,1}\) has order one on each,
and its inverse has order \(-1=-3/b\).

For M6, the complete K94 EU.2-\/-3 formula is
\[
 \begin{gathered}
 U=y-t-x/2,\quad V=y+t+x/2,\quad
 f_2=U/(x-s),\\
 f_3=\frac{s xU}{sV-(s+2t)x},\qquad
 f_6=f_3^2 f_{2,3G},\qquad
 B_{6,1}=f_2^6\,\frac{f_6(p-G)}{f_6(p)},\\
 \left.\frac{B_{6,1}}{t^3}\right|_{C_{\rm id}}
             =\frac{8((1-s)w-s)}{s^3},\qquad
 x|_{C_{\rm id}}=s+w^2-w .
 \end{gathered}                                                  \tag{EF.30b}
\]
The last coefficient is not identically zero for any \(s\ne0\).
For completeness its full factors are
\[
 \begin{aligned}
 f_2(p)&\to x_0/w,&f_3(p)&\to x_0,\\
 f_3(p-G)/t&\to2w^3/[x_0(w-s)],&
 f_{2,3G}(p)&\to x_0(1-s)/[(1-s)w-s],\\
 f_{2,3G}(p-G)/t&\to
                  2(1-s)(w-s)^2/(s^3x_0).
 \end{aligned}                                                  \tag{EF.30c}
\]
Their literal product is EF.30b. The value \(s=1\), which remains
in EF.3, requires specialization of the full map before taking this
limit. Do that in the exact source chart
\[
 u=w-1,\quad v=w+h,\quad x=1+uv,\quad
 t=xh/2,\quad y=x(u+v)/2 .
 \tag{EF.30e}
\]
It satisfies the original equation with \(r=-3/4\). The complete
inverse translation is \(x'=vh,\ y'=t-(v-1/2)x'\). Set \(s=1\)
in the actual tangent at \(3G\) first. Its slope and doubled point are
\[
 \lambda_{3G}=\frac{24t^3+24t^2+3t-1}{2t(4t+3)},\quad
 x(6G)=
 \frac{(4t^2+2t-1)(16t^4+24t^3+16t^2+4t-1)}
      {4t^2(4t+3)^2}.
 \tag{EF.30f}
\]
Substitution of these full functions into the two tangent quotients,
and only then taking \(h\to0\), gives, with \(x_0=w^2-w+1\),
\[
 f_2\longrightarrow x_0/w,\qquad
 \frac{f_3(p-G)}{h f_3(p)}
       \longrightarrow\frac{w^3}{(w-1)x_0},\qquad
 \frac{f_{2,3G}(p-G)}{h f_{2,3G}(p)}
       \longrightarrow-\frac{(w-1)^2}{x_0}.
 \tag{EF.30g}
\]
Since \(t/h\to x_0/2\), their full product gives
\(B_{6,1}/t^3\to-8\) after this specialization. It agrees with EF.30b
but was computed before any generic-\(s\) leading-term cancellation.
The same exact doubling formula gives
\(\tau(6G)/t\to-6\), where \(\tau=-x/y\) at \(O\). On this actual
point/O chart keep the two Cartier equations
\(\tau-\tau(6G)(t)\) and \(\tau\), with both their quotient arrows.
They define the point-line module also at their coincidence. The
Miller identity carries these two modules in the Hom map. No
horizontal divisor is removed by evaluation at its generic point.
After the full order-three factor is removed, all vertical valuations
are zero and the Hom divisor is zero. Normality of the actual smooth
resolved source extends the map and its reciprocal at the remaining
points, including this coincidence.
The section \(6G\) and \(O\) meet the same I2 component, so \(M_6\)
has degree zero on both. The Hom divisor has zero component
intersections. The kernel of the full I2 matrix
\(\left(\begin{smallmatrix}-2&2\\2&-2\end{smallmatrix}\right)\)
is the primitive vector \((1,1)\). Thus the order three in EF.30b
is the order on both components, giving the first M6 entry.

For the root column the exact EX.4 and EX.13 scalings are
\[
 \begin{array}{c|c|c}
 (a,b;m)&B_{a,1,\rm old}&\text{generic central root coefficient}\\ \hline
 (2,3;3l)&z^{-2l}B'_{2,1}&
  B'_{2,1}=f(p)f(p-G'),\quad f=(Y-1)/X\\
 (6,1;3l+2)&z^{-6(l+1)}B'_6&
                         B'_6/z^2\longrightarrow4 .
 \end{array}                                                    \tag{EF.30d}
\]
The first coefficient is nonzero at a generic point of the smooth
root elliptic fibre. The second is the complete chord limit in K96
EX.13, in the line graph EX.10-\/-12. Hence the inverse orders are
\(2l=2m/3\) and \(6(l+1)-2=2m\), respectively. The remaining
root graph conditions are the literal two-section coefficients
EF.32 and the two correction sums EF.33 below, that is K96
EX.11-\/-14; their inverses and characters are carried into EF.35-\/-38.
They are not inferred from this generic order.

The resulting base order of \(H\) at infinity is \(-2j/b\).
The third column is exactly the integer \(B^k\) reindex of EF.24.
At a finite irreducible fibre the numerator and denominator
\(f_a(p-G)/f_a(p)\) have the same vertical order; each other
factor is generically nonzero. This is the component-DVR argument
of K94 EU.2 and K96 EX.1. There are no other fibre divisors.
All graph-correction
monomials are still part of the root modules; EF.30 reads their generic
component orders and does not discard the remaining extension
conditions. The degree of the actual rational divisor is zero:
\(3j/b-2j/b-k=0\), so \(j=bk\). Consequently
\[
 H=c(t^3/q)^k,\qquad H U_G^{bk}=c B_{\rm sel}^{-k},\quad
                                      c\in\mathbb C^*.
 \tag{EF.31}
\]
Since \(B_\lambda=\lambda B_{\rm sel}\) is the actual quotient
deck, EF.31 acts exactly by the physical scalar \(c\lambda^k\).
The equality holds first on the common dense chart and hence on the
entire completed endpoint; EF.7 shows explicitly that this scalar
commutes with all original root, character and annular arrows. This
also retains any additional extension restrictions on the meromorphic
representative in EF.30: the final map is the actual extending scalar,
with the same graph and finite characters.

\par

\medskip\noindent\textbf{EF.7. Both full cap and annular arrows, and the endpoint consequence}

\par\nobreak

Here are the root maps used above, in their original direction. For
\((a,b;m)=(6,1;3l+2)\) put \(k=2l+1\):
\[
 q_\infty=z^m,\quad x_0=z^{2k}X,\quad y_0=z^{3k}Y,\quad
 Y^2=X^3+r z^{2l+2}X^2+z^{l+2}X+z^2.
 \tag{EF.32}
\]
The deck is
\(D(z,X,Y)=(\zeta_mz,\zeta_m^{-2k}X,\zeta_m^{-3k}Y)\).
The three A2 charts of EX.6-\/-8 and their explicit chord inverses give
the actual source resolution and translation. The exact division
section is
\(R=2G'-[m]^{-1}((2m-1)G')\), with \(mR=G'\).
The graph of \(\tau'=z^k\tau_0\), including every ray
\((j,1)\), \(1\le j\le k\), retains separately the section
equations \(\tau'\) and \(\tau'-z^kL(s,z^m)\). Their actual
coefficient ratio is \(u/(u-L)\), or \(1/(1-vL)\) on the inverse
graph chart. It gives \(\beta^*N_6=\pi^*M_6\). The full multiplier
is \(B'_{\rm sel}=h z^{-2}B'_6\),
\(h=1-t_cz^m\), with central coefficient \(4h\).

For \((a,b;m)=(2,3;3l)\) the exact maps are
\[
 q_\infty=z^{3l},\quad x_0=z^{4l}X,\quad y_0=z^{6l}Y,
 \qquad Y^2=X^3+r z^{2l}X^2+z^lX+1.
 \tag{EF.33}
\]
The deck has weights \((1,-4l,-6l)\). Put \(u=z^l\); the
two point-line quotient rays are \((2/3,1/3),(1/3,2/3)\) in
\(\mathbb Z^2+\mathbb Z(2,1)/3\). Their full support difference
is \(\min(v_w,2l v_z)\), at both marked sections. With their
actual monomial generators this gives
\(\beta^*N_2=\pi^*M_2+E_P-E_O\). The source division is the
small section \(R=[m]^{-1}(3G')\). Its translation preserves the
ideals \((w,u^j),(u,w^j)\), \(j=1,2\), and their inverse charts,
so the full line correction is carried. The complete coefficient is
\(B'_{\rm sel}=hB'_{2,3}\); the literal three-orbit Miller product
has central value four. These are the EX.3-\/-5 maps, including the
nontrivial central degree-zero line \(N_2\).

For either branch the actual chord-over-vertical function
\[
 g_{P,R}(p)=\ell_{P,R}(p)/(x(p)-x(P+R))
 \tag{EF.34}
\]
has divisor \((P)+(R)-(P+R)-(O)\) and is the typed lift covering
\(p\mapsto p-R\). The normalized starting lift is
\(F=\tau(P)g_{P,R}\) for \(P=6G'\), and
\(F=\tau(R)g_{P,R}\) for \(P=2G'\). The actual near-O expansion
gives central coefficient one. The formal-coordinate weight and the
chord weight cancel, so it commutes with the full \(D\). The base
unit \(\beta=B'_{\rm sel}/F^m\) has central value four; its
reference-chosen analytic \(m\)-th root is \(D\)-invariant. Thus
\[
 \widetilde F=\beta^{1/m}F,\quad
 A_0=\zeta_m^d\widetilde F^{\,n}D,\quad
 B_\lambda=\lambda B'_{\rm sel},\quad
 A_\lambda=R_\lambda^n A_0,\quad R_\lambda^m=\lambda,
 \tag{EF.35}
\]
and literally \(A_\lambda^m=B_\lambda^n\),
\([A_\lambda,B_\lambda]=1\). These equations retain \(d\),
both signs of \(n\), every source argument of a chord coefficient,
and the selected-height function \(h\).

On the two fixed annular sectors set
\[
 s_\alpha=\log_\alpha q_\infty/(2\pi i),\quad
 z_\alpha=e^{2\pi i s_\alpha/m},\quad
 c_\lambda=n\log R_\lambda/(2\pi i)+d/m,\quad
 f_\alpha=\operatorname{Log}_\alpha\widetilde F.
\]
The full ordinary-to-root map is
\[
 \Phi_\alpha=\operatorname{Exp}
  (n s_\alpha f_\alpha+s_\alpha c_\lambda\delta+h_\alpha)
                                                     I_{z_\alpha}.
 \tag{EF.36}
\]
When \(s_\beta=s_\alpha+j\) and
\(f_\beta=D_*^jf_\alpha+\Pi_{\alpha\beta}\), its actual
reference cochain is
\[
 h_\beta-D_*^jh_\alpha
 =-n s_\alpha\Pi_{\alpha\beta}
       +k_{\alpha\beta}\operatorname{Log}_\beta B_\lambda
       +\Lambda_{\alpha\beta}.
 \tag{EF.37}
\]
The columns \(\Pi,\Lambda\) and integer \(k_{\alpha\beta}\)
are precisely the transported complete reference columns. Subtracting
the target-sheet logarithm \(jn f_\beta+jc_\lambda\delta\)
gives the literal squares
\[
 \Phi_\beta=A_\lambda^jB_\lambda^{k_{\alpha\beta}}\Phi_\alpha,
 \qquad \Phi_\alpha B_\lambda=B_\lambda\Phi_\alpha.
 \tag{EF.38}
\]
Their inverses use the same columns with inverse group elements. The
full source displacement is the signed \(-n s_\alpha\) Abel
integral of \(R\), and the scalar is
\(R_\lambda^{n s_\alpha}e^{2\pi i d s_\alpha/m}\).
This is the complete cap-period extension used in EF.3 and EF.6.

Physical multiplication by \(c\) commutes with EF.34-\/-38 and acts
on each cusp character by \(Z^uV^vq^h\mapsto c^vZ^uV^vq^h\).
It therefore extends on the same root and fan with the prescribed
reference. A path in \(\mathbb C^*\) from one to \(c\) is an
actual isotopy of these completed maps. Conversely EF.2-\/-6 show every
holomorphic automorphism is one of them. Its kernel is trivial: on a
regular source fibre an equality with a suspension power would give
an integral multiple of \(bG\) equal to zero as a global source
section, contradicted by EF.21; then the scalar itself is one.

For two actual holomorphic endpoint identifications of one and the
same cut family, their ratio is such an automorphism of the common
completed endpoint. The physical return uses the inverse endpoint;
the scalar path therefore supplies an actual isotopy between those
inverse-endpoint returns, with the specified source, fan and cap
markings carried. This conclusion concerns endpoints that have been
constructed. It assigns neither a new shorter scalar endpoint nor a
return to an unconstructed source loop.

\par

\medskip\noindent\textbf{EF.8. Exact verification and scope ledger}

\par\nobreak

The exact algebra is evaluated by factoring rational functions over the
displayed source equation, using symbolic matrix multiplication for the full
central cochain, and computing the determinant of the ten-class intersection
matrix. In particular the $M_6$ coefficient at $s=1$ is formed in the original
$x$ coordinate before its central limit is taken. The program
\path{code/current_audits/audit_corrections/code/independent_IV_endpoint_checks.py} and output \path{code/current_audits/audit_corrections/independent_IV_endpoint_checks.json} contain 25,796 exact
checks across \cref{night:P3K.109,night:P3K.110,night:P3K.113,night:P3K.114}.
The checks used here include the discriminant and marked-source inverse, the
$2G/3G$ chord identities, the full specialized $M_6$ coefficient, the section
lattice, and the symbolic period and dual equations. The Lie extension,
proper fixed-curve closure, covering and annular maps are proved above.

The source, root and reference inputs are the explicit constructions in
\cref{night:P3K.96,night:P3K.97} and the topology maps of
\cref{night:P3F.35,night:P3F.46}. The source equations, section and divisor
orders used throughout are those displayed in EF.1 and EF.18.





% End expanded verified cp36 input

\end{proof}


\subsection{All holomorphic scalar endpoints on the generic fixed E6 source}
